%%%%%%%%%%%%%%%%%%%%%%%%%%%%%%%%%%%%%%%%%%%%%%%%%%%%%
%			HLAVIČKA								%
%%%%%%%%%%%%%%%%%%%%%%%%%%%%%%%%%%%%%%%%%%%%%%%%%%%%%
\documentclass[openany,12pt]{memoir}

\usepackage[utf8]{inputenc} 
\usepackage[czech]{babel}
\usepackage[T1]{fontenc}
\usepackage[top=1.5cm, bottom=2cm, left=2cm, right=2cm]{geometry}  % --> NASTAVENÍ OKRAJŮ
\usepackage{fancyhdr}
\usepackage{graphicx}
\usepackage{lmodern}
\usepackage{xwatermark}
\usepackage{xcolor}
\usepackage{changepage}
\usepackage{newunicodechar}



%%%%%% Package na zpěvník
\usepackage[full]{leadsheets}%http://mirrors.nic.cz/tex-archive/macros/latex/contrib/leadsheets/leadsheets_en.pdf   --> dokumentace	
\definesongtitletemplate{empty}{} 
\setchords{
format = \bfseries,   %tučné akordy
minor = {mi},% 
input-notation = {german},%
output-notation = {german}%
}
\definesongtitletemplate{empty}{} 

\newlength{\drop}
\newwatermark[allpages,color=red!50,angle=0,scale=2, xpos=0,ypos=0]{\includegraphics[width=5cm]{obr/pozadi.jpg}} %--> dvojka na pozadí


%%%% Vlastní příkazy
\newcounter{Slokočet}   %Automatické číslování slok
\newcommand{\mezera}{\vspace*{0.5cm}}   %Horizontální odsazení slok
\newcommand{\stred}{5.2cm}   %%% Na zarovnání slok doprostřed, pozn. automatičtější zarovnávání na střed nejde
\newcommand{\refren}{\mezera \noindent \textbf{R:} } %refrén
\newcommand{\sloka}{\mezera \noindent \addtocounter{Slokočet}{1} \arabic{Slokočet}. } 	%sloka, která se automaticky čísluje
\newcommand{\ssloka}{\mezera \noindent}  % vlastní číslo sloky
\newcommand{\carka}{,\:}
\newcommand{\m}[1]{\color{white}{#1}}  %Pro akordy

%%%%%%%%%%%%%%%%%%%%%%%%%%%%%%%%%%%%%%%%%%%%%%%%%%%%%%
%			NÁVOD									 %
%%%%%%%%%%%%%%%%%%%%%%%%%%%%%%%%%%%%%%%%%%%%%%%%%%%%%%
%1. Věci v hlavičce IGNOROVAT
%2. Píseň psát do prostoru mezi \begin{song} a \end{song}
%3. další řádek se značí dvěma odsazeníma (= dvakrát stisknout enter)
%4. \refren vždy na začátku refrenu a \sloka na začátek sloky (automaticky se čísluje)
%  \ = alt gr + q ; [/] = alt gr f/g ; {/} = alt gr + b/n; ^ = alt gr + 3 + mezera
%Cokoliv napíšete do ^{  } se bude brát jako akord
%když se toto bude dotýkat nějakého slova (nebude mezi tím a slovem mezera)
%tak se akordy zjeví nad slovem, ale pište to před slovo
%Když se to nedotýká slova, tak akord lítá ve vzduchu a vytiskne se větší mezera
%První možnost je asi preferovanější
%5. Akordy stačí psát jen do první sloky, když se nezmění -- kytaristi to zvládnou



%<++>
%\usepackage{subfiles}
%</++>

\begin{document}

\pagestyle{empty}
%\addtocontents{toc}{\protect\thispagestyle{empty}}
%\tableofcontents \thispagestyle{empty}\newpage
\newgeometry{top=1.5cm, bottom = 0cm, left = 2cm, right = 2cm}




%Pro kompilaci je důležitý, aby ten input a \end{document} zůstali na posledních
%dvou řádcích!!!!
%\documentclass[../main.tex]{subfiles}

\begin{song}{title=\centering Lemon Tree \\\normalsize Fool's Garden  \vspace*{-0.3cm}}  %% sem se napíše jméno songu a autor
\moveright 1cm \vbox{      %Varianta č. 1  ---> Jeden sloupec zarovnaný na střed	
\begin{minipage}[t]{0.48\textwidth}\setlength{\parindent}{0.45cm}  %Varianta č. 2 --> Dva sloupce

\sloka
I'm ^{Emi{\color{white}\_}}sittin' here in the ^{Hmi}boring room,

It's ^{Emi}just another rainy Sunday ^{Hmi}afternoon.

I'm ^{Emi{\color{white}\_\_}}wasting my time I got ^{Hmi}nothing to do,

I'm ^{Emi{\color{white}\_}}hangin' aroud I got ^{Hmi}waitin' for you.

But ^{Ami{\color{white}\_\_}}nothing ever happens^{H}

And I ^{Emi{\color{white}\_\_}}wonder.

\sloka
I'm drivin' aroud in my car,

I'm drivin' too fast, I'm drivin' too far,

I'd like to change my point of view,

I feel so lonely I'm waiting for you.

But nothing ever happens.

\refren
^{G}I wonder how, ^{D\,\,\,\,\,\,\,\,\,}wonder why

^{Emi{\color{white}\_\_}\,\,}Yesterday you told me 'bout the ^{Hmi}blue

 ^{\phantom{.}}blue sky

And ^{C}all that I can ^{D}see

Is just a yellow ^{G\,\,\,\,\,\,}lemon tree.^{D} 

I'm ^{G{\color{white}\_\_}}turnin' my head ^{D}up and down.

^{Emi}I'm turnin', turnin', turnin', ^{Hmi}turnin',

 ^{\phantom{.}}turnin' around

And ^{C}all that I can ^{D}see

Is just a yellow ^{G{\color{white}\_\_}}lemon tree.^{D} 

\sloka
^{4(Emi Hmi) Ami H Emi}Sing\,dip,\,ta\,da\,da\,dap\,di\,dap\,dam\elipsa.\elipsa.\elipsa.

\end{minipage}\begin{minipage}[t]{0.48\textwidth}\setlength{\parindent}{0.45cm}  % V případě varianty č.2 jde odsud text do pravé části
\vspace*{0.58cm}
\sloka
I'm sittin' here, I miss the power,

I'd like to go out, takin' a shower.

But there's a heavy cloud inside my head,

I feel so tired put myself into bed

Where nothing ever happens

And I wonder.

\sloka
^{H{\color{white}\_\_\_}}Isolation ^{Emi}is not good for me.

^{D{\color{white}\_\_\_}}Isolation, ^{G}I don't want to

^{H}sit on a lemon tree.

\sloka
I'm steppin' around in a desert of joy,

Baby, anyhow I get another toy

And everything will happen

And you'll wonder.

\refren
%^{G}I wonder how\elipsa\dots \\

\refren
%^{G}I wonder how\elipsa\dots \\


\dots\,and all that I can see --

and all that I can see --

and all that I can see 

Is just a yellow lemon tree.

\end{minipage}
}
\setcounter{Slokočet}{0}
\end{song}
\end{document}
